% ============================================================
% chapter-state-of-art.tex
% Capítulo 3: Estado del Arte
% ============================================================

\chapter{Estado del Arte}
\label{ch:state-of-art}

\section{Introducción}
[Presentar el objetivo de la revisión.]

\section{Objetivos de revisión}
[Describir qué se busca responder con la revisión.]

\section{Preguntas de revisión}
\begin{itemize}
  \item P1. ¿[Texto de la pregunta 1]?
  \item P2. ¿[Texto de la pregunta 2]?
\end{itemize}

\section{Estrategia de búsqueda}

\subsection{Motores de búsqueda}
[Incluir al menos dos motores de búsqueda.]

\subsection{Cadenas de búsqueda}
[Indicar las cadenas usadas en cada motor.]

\subsection{Documentos encontrados}

\begin{table}[H]
  \centering
  \caption{Documentos encontrados por pregunta}
  \label{tab:documents}
  \begin{tabular}{|l|c|c|}
    \hline
    \textbf{Motor de búsqueda} & \textbf{P1} & \textbf{P2} \\
    \hline
    [Motor 1] & [n] & [n] \\
    \hline
    [Motor 2] & [n] & [n] \\
    \hline
  \end{tabular}
\end{table}

\subsection{Criterios de inclusión/exclusión}
[Describir los criterios aplicados.]

\section{Formulario de extracción de datos}
[Presentar el formulario usado para extraer información
de los estudios primarios.]

\section{Resultados de la revisión}

\subsection{Respuesta a P1}
[Elaborar respuesta basada en el formulario de extracción.
No resumir estudios primarios directamente.]

\subsection{Respuesta a P2}
[Ídem.]

\section{Conclusiones}
[Conclusiones de la revisión del estado del arte.]
